% This is text associated with the open textbook "Open Electromagnetics"
% See immediately below for author, date
% This work is licensed under the Creative Commons Attribution-ShareAlike 4.0 International License. (CC BY-SA 4.0) 
% To view a copy of the license, visit http://creativecommons.org/licenses/by-sa/4.0/

%ORB TITLE Impedance of the Electrically-Short Dipole
%ORB AUTHOR 0        # S.W. Ellingson, Virginia Tech (ellingson@vt.edu)
%ORB DATE 20191127   # yyyymmdd
%ORB ABSTRACT 

\label{m0204_The_Electrically-Short_Dipole} % <-- Must match name of this file and module directory name.  Don't mess with this!
\index{antenna!electrically-short dipole (ESD)|(}

In this section, we determine the impedance of the electrically-short dipole (ESD) antenna.  
The physical theory of operation for the ESD is introduced in Section~\ref{m0198_Radiation_from_an_Electrically-Short_Dipole}, and additional relevant aspects of the ESD antenna are addressed in Sections~\ref{m0201_How_Antennas_Radiate}--\ref{m0210_Reactance_of_the_Electrically-Short_Dipole}. 
The concept of antenna impedance is addressed in Section~\ref{m0202_Equivalent_Circuit_Model_for_Transmission}.
Review of those sections is suggested before attempting this section.

The impedance of any antenna may be expressed as
%
\begin{equation}
Z_A = R_{rad} + R_{loss} +jX_A
\end{equation}
%
where the real-valued quantities $R_{rad}$, 
\index{radiation resistance}
$R_{loss}$, and $X_A$ represent the radiation resistance, loss resistance, and reactance of the antenna, respectively.

{\bf Radiation resistance.}
The radiation resistance of any antenna can be expressed as:
%
\begin{equation}
R_{rad} = 2P_{rad} \left|I_0\right|^{-2}
\end{equation}
%
where $\left|I_0\right|$ is the magnitude of the current at the antenna terminals, and $P_{rad}$ is the resulting total power radiated. 
For an ESD (Section~\ref{m0207_Total_Power_Radiated_by_an_Electrically-Short_Dipole}):
%
\begin{equation}
P_{rad} \approx \eta \frac{\left|I_0\right|^2\left(\beta L\right)^2}{48\pi}
\label{m0204_ePT}
\end{equation}
%
so
%
\begin{equation}
R_{rad} \approx \eta \frac{\left(\beta L\right)^2}{24\pi}
\label{m0204_eRrad}
\end{equation}
%
It is useful to have an alternative form of this expression in terms of wavelength $\lambda$.
This is derived as follows.
First, note:
%
\begin{equation}
\beta L = \frac{2\pi}{\lambda} L = 2\pi\frac{L}{\lambda}
\end{equation}
%
where $L/\lambda$ is the antenna length in units of wavelength.
Substituting this expression into Equation~\ref{m0204_eRrad}:
%
\begin{flalign}
R_{rad} &\approx \eta \left(2\pi\frac{L}{\lambda}\right)^2 \frac{1}{24\pi} \nonumber \\
        &\approx \eta \frac{\pi}{6} \left(\frac{L}{\lambda}\right)^2
\end{flalign}
%
Assuming free-space conditions, $\eta\cong 376.7~\Omega$, which is $\approx 120\pi~\Omega$. Subsequently,
%
\begin{equation}
\boxed{
R_{rad} \approx 20\pi^2 \left(\frac{L}{\lambda}\right)^2
}
\label{m0204_eRrad2}
\end{equation}
%
This remarkably simple expression indicates that the radiation resistance of an ESD is very small (since $L\ll\lambda$ for an ESD), but increases as the square of the length.

At this point, a warning is in order.  
The radiation resistance of an ``electrically-short dipole'' is sometimes said to be $80\pi^2 \left(L/\lambda\right)^2$; i.e., 4 times the right side of Equation~\ref{m0204_eRrad2}.  
This higher value is not for a \emph{physically-realizable} ESD, but rather for the \emph{Hertzian dipole} (sometimes also referred to as an ``ideal dipole'').
\index{Hertzian dipole}
The Hertzian dipole is an electrically-short dipole with a current distribution that has uniform magnitude over the length of the dipole.\footnote{See Section~\ref{m0197_Radiation_from_a_Hertzian_Dipole} for additional information about the Hertzian dipole.}
The Hertzian dipole is quite difficult to realize in practice, and nearly all practical electrically-short dipoles exhibit a current distribution that is closer to that of the ESD.
The ESD has a current distribution which is maximum in the center and goes approximately linearly to zero at the ends.
The factor-of-4 difference in the radiation resistance of the Hertzian dipole relative to the (practical) ESD is a straightforward consequence of the difference in the current distributions.

{\bf Loss resistance.}  The loss resistance of the ESD is derived in Section~\ref{m0208_Total_Power_Dissipated_by_an_Electrically-Short_Dipole} and is found to be
%
\begin{equation}
R_{loss} \approx \frac{L}{6a} \sqrt{\frac{\mu f}{\pi \sigma}}  
\label{m0204_eRloss}
\end{equation}
%
where $a$, $\mu$, and $\sigma$ are the radius, permeability, and conductivity of the wire comprising the antenna, and $f$ is frequency.

{\bf Reactance.} The reactance of the ESD is addressed in Section~\ref{m0210_Reactance_of_the_Electrically-Short_Dipole}.  A suitable expression for this reactance is (see, e.g., Johnson (1993) in ``Additional References'' at the end of this section):
%
\begin{equation}
X_A \approx -\frac{120~\Omega}{\pi L/\lambda} \left[ \ln{\left(\frac{L}{2a}\right)}-1 \right]
\label{m0204_eXA}
\end{equation}
%
where $a\ll L$ is assumed.


%%%%%%%%%%%%%%%%%%%%%%%%%%%%%%%%%%%%%%%%%%%%%%%
%ORB EXAMPLE
~\\ \begin{example} 
\begin{mdframed}[backgroundcolor=brown!20,linecolor=brown!20]\raggedright
Impedance of an ESD.
\\~\\
A thin, straight dipole antenna operates at 30~MHz in free space.  
The length and radius of the dipole are 1~m and 1~mm respectively. 
The dipole is comprised of aluminum having conductivity $\approx 3.7 \times 10^7$~S/m and $\mu\approx\mu_0$.
What is the impedance and radiation efficiency of this antenna?

{\bf Solution.}
The free-space wavelength at $f=30$~MHz is $\lambda=c/f\cong10$~m.
Therefore, $L\cong 0.1\lambda$, and this is an ESD.
Since this is an ESD, we may compute the radiation resistance using Equation~\ref{m0204_eRrad2}, yielding $R_{rad} \approx 1.97~\Omega$.
The radius $a=1$~mm and $\sigma\approx 3.7 \times 10^7$~S/m; thus, we find that the loss resistance $R_{loss} \approx 94.9~\mbox{m}\Omega$ using Equation~\ref{m0204_eRloss}.
Using Equation~\ref{m0204_eXA}, the reactance is found to be $X_A \approx -1991.8~\Omega$.
The impedance of this antenna is therefore
%
\begin{flalign}
Z_A &= R_{rad} + R_{loss} +jX_A \nonumber \\
    &\approx \underline{2.1-j1991.8~\Omega}
\end{flalign}
%
and the radiation efficiency of this antenna is
%
\begin{equation}
e_{rad} = \frac{R_{rad}}{R_{rad}+R_{loss}} \approx \underline{95.4\%}
\end{equation}
%
\end{mdframed}
% note figures will need to go between "\end{mdframed}" and "\end{example}"
\end{example}
%%%%%%%%%%%%%%%%%%%%%%%%%%%%%%%%%%%%%%%%%%%%%%%

In the preceding example, the radiation efficiency is a respectable 95.4\%. 
However, the real part of the impedance of the ESD is much less than the characteristic impedance of typical transmission lines (typically 10s to 100s of ohms).
Another problem is that the reactance of the ESD is very large.
Some form of impedance matching is required to efficiently transfer power from a transmitter or transmission line to this form of antenna.
Failure to do so will result in reflection of a large fraction of the power incident on antenna terminals.
A common solution is to insert series inductance to reduce -- nominally, cancel -- the negative reactance of the ESD.
In the preceding example, about $10~\mu$H would be needed.
The remaining mismatch in the real-valued components of impedance is then relatively easy to mitigate using common ``real-to-real'' impedance matching techniques. 

%%%%%%%%%%%%%%%%%%%%%

{\bf Additional Reading:}
\begin{itemize}
\item R.C.\ Johnson (Ed.), \emph{Antenna Systems Handbook} (Ch.\ 4), McGraw-Hill, 1993.
\end{itemize}

%%%%%%%%%%%%%%%%%%%%%%

\index{antenna!electrically-short dipole (ESD)|)}

